\section*{Chapitre \ref{ch:qm:finite-difference-methods} --- Méthodes des différences finies}

\begin{solution}{exo:qm:finite-difference-methods:1}
$\Delta\tau \le \Delta x^2/\sigma^2 = 6.25 \times 10^{-6}/0.04 = 1.5625 \times 10^{-4}$ années, soit 6\,400 pas pour un an.
\end{solution}

\begin{solution}{exo:qm:finite-difference-methods:2}
$4.3572 + (4.3572 - 4.3558)/3 = 4.35767$.
\end{solution}

\begin{solution}{exo:qm:finite-difference-methods:3}
$u(x \pm h) = u \pm hu_x + \frac{h^2}{2}u_{xx} \pm \frac{h^3}{6}u_{xxx} + \frac{h^4}{24}u_{xxxx} + O(h^5)$~; en additionnant et en retranchant $2u$, puis en divisant par $h^2$, il reste $u_{xx} + \frac{h^2}{12}u_{xxxx} + O(h^4)$.
\end{solution}

\begin{solution}{exo:qm:finite-difference-methods:4}
$g(\pi) = (1 - 64)/(1 + 64) = -0.969$~; $|g|^n = 0.01$ exige $n = \ln100/(-\ln0.969) \approx 148$ pas, bien plus que les 25 de la simulation entière. Un demi-pas implicite en $\lambda/2 = 16$ a pour facteur $1/(1 + 64) =
0.015$ en $\xi = \pi$~: un seul demi-pas amortit ce mode d'un facteur 65.
\end{solution}

\begin{solution}{exo:qm:finite-difference-methods:5}
Les coefficients valent $a/h^2 \mp b/(2h)$~; tous deux sont positifs ou nuls si et seulement si $h \le 2a/|b|$, c'est-à-dire $|b|h/(2a) \le 1$.
\end{solution}

\begin{solution}{exo:qm:finite-difference-methods:6}
Loin sous le strike, le call vaut presque exactement zéro, donc $u = 0$ est presque la valeur vraie et son petit écart est amorti par la diffusion avant d'atteindre le strike. Loin au-dessus, le call vaut
environ $S - Ke^{-r\tau}$, un nombre élevé~: imposer zéro à cet endroit introduit une grande erreur qui diffuse vers l'intérieur et, sur la durée de vie de l'option, atteint le strike.
\end{solution}

\begin{solution}{exo:qm:finite-difference-methods:7}
\texttt{digital\_placement()}~: sur un nœud $9.9 \times 10^{-3}$, $4.9 \times 10^{-3}$, $2.5 \times 10^{-3}$, $1.2 \times 10^{-3}$ (rapports 2.0, premier ordre)~; entre les nœuds $6.6 \times 10^{-6}$, $1.5 \times 10^{-6}$,
$3.8 \times 10^{-7}$, $9.5 \times 10^{-8}$ (rapports 4.4, 4.0, 4.0, second ordre).
\end{solution}

\begin{solution}{exo:qm:finite-difference-methods:8}
Un prix convergé ne renseigne en rien sur ses dérivées~: une oscillation d'un nœud à l'autre d'amplitude $\epsilon$ dans les valeurs devient $\epsilon/h$ dans le delta et $\epsilon/h^2$ dans le gamma. Dans la simulation du chapitre, le prix est à moins de
0.018 et le gamma erroné de vingt fois sa propre taille. Les grecques ont besoin de leur propre test de convergence.
\end{solution}

\begin{solution}{pb:qm:finite-difference-methods:1}
\textbf{1.} $\lambda = a\Delta\tau/\Delta x^2 = 0.02 \times 0.01/0.0025^2 = 32$~; le \omterm{def:qm:finite-difference-methods:theta}{schéma explicite} exige $\lambda \le \frac12$.
\textbf{2.} À $\lambda = \frac12$~: pour 0.489 (409 pas sur 200 nœuds), le prix vaut 4.3576~; pour 0.509 (393 pas), la plus grande valeur de la grille vaut 2\,413.
\textbf{3.} $(1 - 2\lambda)/(1 + 2\lambda) = -63/65 = -0.969$.
\textbf{4.} Le prix est à moins de 0.018 de 4.3576~; le gamma oscille entre $-0.81$ et $0.55$ autour de 0.040 (erreur maximale 0.84).
\textbf{5.} Elle double~: 0.46, 0.84, 1.69 et 3.38 pour 200 à 1\,600 nœuds, car $\lambda$ double.
\textbf{6.} Erreur maximale sur le gamma $4.7 \times 10^{-4}$ pour 400 nœuds, divisée par deux à chaque raffinement~: premier ordre.
\textbf{7.} $1.76 \times 10^{-5}$, $4.35 \times 10^{-6}$, $1.07 \times 10^{-6}$~: rapports 4.05 et 4.07, second ordre.
\textbf{8.} Premier ordre~: 0.041, 0.018, 0.0092, 0.0046.
\textbf{9.} Une discontinuité échantillonnée à un nœud (valeur $\frac12$ en ce point) laisse une erreur au niveau du strike, proportionnelle au pas en espace, comme le montrent les rapports mesurés égaux à 2~; la perte vient de la manière dont
les données discontinues sont échantillonnées, et non du schéma en temps, puisque des pas de temps plus petits ne la modifient pas.
\textbf{10.} Erreur $6.6 \times 10^{-6}$ pour 200 nœuds, diminuant d'un facteur 4 à chaque fois que le pas est divisé par deux.
\textbf{11.} Erreur environ 2.9 fois plus petite à nombre de nœuds égal (0.0079 contre 0.0225 pour 50 nœuds).
\textbf{12.} De $1.8 \times 10^{-3}$ et $4.6 \times 10^{-4}$ à $3.2 \times 10^{-7}$. Sans le démarrage, l'erreur sur le prix est de premier ordre, et l'extrapolation avec $p = 2$ élimine le terme fautif.
\textbf{13.} Lorsque $|b|\Delta x/(2a) > 1$~; ici 2.49, et le digital atteint 1.030 (au-dessus du paiement actualisé 0.951) avec une variation totale de 1.23~; le \omterm{def:qm:finite-difference-methods:m}{schéma décentré amont} rétablit 0.951 et la monotonie.
\textbf{14.} Gamma nul et bords de Dirichlet donnent, au strike, des prix qui diffèrent de $6 \times 10^{-8}$ avec une grille à cinq écarts-types.
\textbf{15.} Erreurs de 0.045, 0.026 et 0.014 par rapport à la valeur de Margrabe 10.5243 sur $40^2$, $80^2$, $160^2$ nœuds~: premier ordre, dû au coude le long de la diagonale entre les nœuds.
\textbf{16.} Un test de raffinement sur les grecques (le gamma doit voir son erreur diminuer d'un facteur 4 à chaque division par deux, et non croître), ou un décompte des changements de signe de la différence seconde près du strike.
\textbf{17.} La stabilité borne la croissance de chaque mode, mais non son signe~: le \omterm{def:qm:finite-difference-methods:theta}{schéma de Crank--Nicolson} maintient en vie les hautes fréquences avec un signe alterné, et la différentiation discrète deux fois de suite les amplifie d'un facteur $1/h^2$.
\textbf{18.} Démarrage implicite sur chaque donnée non lisse, strikes et barrières entre les nœuds, ordres de convergence mesurés pour les prix et les grecques, vérification de la \omterm{def:qm:finite-difference-methods:m}{M-matrice}, et extrapolation seulement à un ordre mesuré.
\textbf{19.} Phénomène nommé~: \emph{le gamma en dents de scie}~: quatre demi-pas implicites suppriment l'oscillation et rétablissent le second ordre, l'erreur sur le gamma diminuant d'un facteur 4.05 puis 4.07 à chaque division par deux, là où celle du Crank--Nicolson simple
double, et deux demi-pas donnent un premier ordre.
\textbf{20.} Ce que le schéma fait à chaque fréquence à chaque pas, ce que voient les grecques.
\end{solution}

\begin{solution}{iq:qm:finite-difference-methods:1}
$\Delta t \le \Delta x^2/(2a)$ pour $u_t = au_{xx}$ ($\Delta t \le \Delta x^2/\sigma^2$ en log-prix). Diviser $\Delta x$ par deux divise le pas en temps par quatre, de sorte qu'une grille fine rend le \omterm{def:qm:finite-difference-methods:theta}{schéma explicite} très coûteux, ce qui motive les
\omterm{def:qm:finite-difference-methods:theta}{schémas implicites}.
\iqlookfor{La mise à l'échelle en $\Delta x^2$ et son coût.}
\end{solution}

\begin{solution}{iq:qm:finite-difference-methods:2}
Son facteur d'amplification tend vers $-1$ aux hautes fréquences lorsque $\Delta t/\Delta x^2$ est grand, de sorte que le contenu haute fréquence d'un paiement à coude survit avec un signe alterné et apparaît dans la
différence seconde. Remède~: quelques pas implicites (demi-pas) au départ (\omterm{def:qm:finite-difference-methods:rannacher}{démarrage de Rannacher}), lissage ou placement du paiement, ou un schéma L-stable.
\iqlookfor{Le facteur d'amplification proche de $-1$ et le \omterm{def:qm:finite-difference-methods:rannacher}{démarrage de Rannacher}.}
\end{solution}

\begin{solution}{iq:qm:finite-difference-methods:3}
À mi-chemin entre les nœuds (ou avec le paiement moyenné sur chaque cellule) pour les strikes~; exactement sur un nœud pour une barrière surveillée en continu, où la condition aux limites est imposée. L'erreur liée à une
discontinuité située dans une cellule dépend de sa position dans celle-ci et détruit l'ordre de convergence.
\iqlookfor{Le placement détermine l'ordre.}
\end{solution}

\begin{solution}{iq:qm:finite-difference-methods:4}
En dimension un, le système tridiagonal se résout par l'algorithme de Thomas en $O(n)$. En dimension deux, la méthode ADI résout un système tridiagonal par ligne de grille et par direction, soit $O(n^2)$ par pas, au lieu
d'une résolution creuse du système complet.
\iqlookfor{Thomas et ADI.}
\end{solution}

\begin{solution}{iq:qm:finite-difference-methods:5}
Raffiner l'espace et le temps ensemble par facteur deux sur un problème à solution connue (ou par comparaison avec une référence bien plus fine), et vérifier que le rapport d'erreur tend vers $2^p$ pour les prix et les grecques séparément,
en des points qui restent sur la grille.
\iqlookfor{Rapports d'erreur à chaque division par deux, grecques comprises.}
\end{solution}

\begin{solution}{iq:qm:finite-difference-methods:6}
Le schéma n'est pas monotone~: les différences premières centrées avec une dérive grande devant la diffusion (ou un stencil à dérivées croisées avec forte corrélation) donnent des coefficients hors diagonale négatifs, et
le \omterm{def:qm:finite-difference-methods:theta}{schéma de Crank--Nicolson} n'amortit pas les oscillations qui en résultent. Recourir au \omterm{def:qm:finite-difference-methods:m}{schéma décentré amont}, raffiner là où la dérive domine, employer un stencil à dérivées croisées préservant la positivité, et démarrer par des pas implicites.
\iqlookfor{Perte de la propriété de \omterm{def:qm:finite-difference-methods:m}{M-matrice}.}
\end{solution}
